% 本文件由 assemble.py 按 _work/plans/ch99.json 逐字组装，请勿手改（改 plan 或 blocks 后重新生成）。
\chapter{其他学科与未归类笔记}
\label{ch:99}

\section{页边零星·化学}
% ---- block 081-01 (原稿第81页) kind=misc ----
% 页面上缘空白处(页眉“No.”旁)有零星字迹，疑为相邻页内容(化学·气体密度)，照录
$\rho=1.429\unit{g\cdot L^{-1}}$

% ---- block 083-01 (原稿第83页) kind=misc ----
% 页面上缘右侧页眉“No.”旁有零星字迹(与p081上缘相同，疑为化学·气体密度)，照录
$\rho=1.429\unit{g\cdot L^{-1}}$

% ---- block 085-01 (原稿第85页) kind=misc ----
% 页面上缘右侧页眉“No.”旁有零星字迹(与p081、p083上缘相同，疑为化学·气体密度)，照录
$\rho=1.429\unit{g\cdot L^{-1}}$


\section{化学·吸附态相对能量(非物理内容)}
% ---- block 196-02 (原稿第196页) kind=note ----
% 与物理无关的化学笔记（含吸附态记号 $*$）；“相对能量”四字为写在 CH$_3$OH$^*$ 上方的小字
% 此行左下方(y约0.355-0.37,x约0.02-0.06)另有一处小记号(一道竖线加两三道斜线，含义不明，疑似涂鸦)
CH$_3$OH$^{*}$\textsuperscript{相对能量}\ 不是甲醇的能量\qquad CO$^{*}$\ 能量不是CO能量


\section{生物·标志重捕法与稀释涂布平板计数}
% ---- block 199-03 (原稿第199页) kind=note ----
% 与物理无关的生物笔记；后三行的“$\rho$小”被同一个椭圆圈起，首行“$\rho$大”下有一道横线
\begin{itemize}
\item 标志重捕法\ 标志物脱落\quad $\rho$大
\item 计数酵母菌时未稀释\quad $\rho$小
\item 在稀疏区取样\quad $\rho$小
\item 用稀释涂布平板计数大肠杆菌\quad $\rho$小
\end{itemize}


\section{化学·有机反应(内成环加成)}
% ---- block 199-04 (原稿第199页) kind=note ----
% 与物理无关的化学笔记；红笔下划线共三处（“内成环”、“也”、“加成反应”，文字本身为黑色）
分子\red{\uline{内成环}}\red{\uline{也}}叫\red{\uline{加成反应}}，（啥也没加）。


\section{化学·氢化物与硅硫硼化合物}
% ---- block 199-05 (原稿第199页) kind=note ----
% 与物理无关的化学笔记。红笔：首句“氢化”下、方框内“最简单氢化物”下、两个圆圈内文字下各有一道红色下划线；方框内压着一个五角星
\red{\uline{氢化}}物注意\ \fbox{最简单氢化物 ★}\ \ \fbox{CH$_4$}\ 与\ \fbox{\unclear{沥青}}\ 都是氢化物

\notefig[0.6]{figures/p199_d.png}
% 结构式转录：(CH$_3$)$_3$ *Si$-$S$-$B，B 上分两支：上支 $-$S$-$Si(CH$_3$)$_3$；下支 S \struck{CH$_3$} 另起 Si$-$(CH$_3$)$_3$

